% Source: https://artofproblemsolving.com/community/c3h2620959_geometry_problem
% Stable identity: aops-topic-2620959-q01
Diagonal $BD$ splits the quadrilateral into triangles $ABD$ and $BCD$. By the triangle inequality:
\begin{itemize}
\item in triangle $ABD$, $BD<AB+DA=5+9=14$;
\item in triangle $BCD$, $BD+CD>BC$, i.e. $BD>17-5=12$.
\end{itemize}
So $12<BD<14$, and the integer is $BD=\boxed{13}$. This value is possible: triangles with sides $5,9,13$ and $5,13,17$ exist, and placing them on opposite sides of a common side $BD=13$ gives the quadrilateral.

\begin{asy}
usepackage("amsmath");
size(6.5cm);
defaultpen(fontsize(10pt));
real bx=113/26,cx=433/26;
pair B=(0,0),D=(13,0),A=(bx,sqrt(25-bx^2)),C=(cx,-sqrt(289-cx^2));
draw(A--B--C--D--cycle,linewidth(0.8));draw(B--D,dashed);
label("$A$",A,N);label("$B$",B,W);label("$C$",C,SE);label("$D$",D,NE);
label("$5$",(A+B)/2,NW);label("$9$",(A+D)/2,NE);label("$17$",(B+C)/2,SW);label("$5$",(C+D)/2,E);label("$13$",(B+D)/2+(1.5,0),N);
\end{asy}
